\begin{proof}[Proof of \cref{obj:lem:positive-family-certificate}]
\begin{enumerate}
\item \textbf{Frame identities.}
Write the active index set from \cref{obj:def:cosine-family} as
\[
\begin{split}
\mathcal J=\biggl\{j\in\mathbb Z:\;&
\left\lfloor\frac{x_0-h_{\mathrm L}}{\delta_{\mathrm L}}\right\rfloor-1
\le j\le
\left\lceil\frac{x_0+h_{\mathrm L}}{\delta_{\mathrm L}}\right\rceil+1,\\
&j\delta_{\mathrm L}-\delta_{\mathrm L}<x_0+h_{\mathrm L},
\quad
x_0-h_{\mathrm L}<j\delta_{\mathrm L}+\delta_{\mathrm L}
\biggr\}.
\end{split}
\]
For \(x\in[0,1]\), \(j\in\mathbb Z\), and
\(\lambda\in\{0,1\}^{\mathcal J}\), define
\[
g(x)=
\begin{cases}
\cos^2\!\left(\dfrac{\pi(x-x_0)}{2h_{\mathrm L}}\right),
 & |x-x_0|\le h_{\mathrm L},\\
0,& |x-x_0|>h_{\mathrm L},
\end{cases}
\qquad
\phi_j(x)=
\begin{cases}
\cos\!\left(\dfrac{\pi(x-j\delta_{\mathrm L})}{2\delta_{\mathrm L}}\right),
 & |x-j\delta_{\mathrm L}|\le\delta_{\mathrm L},\\
0,& |x-j\delta_{\mathrm L}|>\delta_{\mathrm L},
\end{cases}
\]
and
\[
\sigma_{\lambda,j}=2\lambda_j-1\quad(j\in\mathcal J),
\qquad
F_\lambda(x)=\sum_{j\in\mathcal J}\sigma_{\lambda,j}\phi_j(x).
\]
Thus
\[
q(x)=g(x)^2,\qquad S_\lambda(x)=g(x)F_\lambda(x),
\qquad
0\le g(x)\le1,\quad 0\le q(x)\le1,\quad |\phi_j(x)|\le1.
\]
The frame functions vanish at their support endpoints.

To establish the partition identity, put
\[
j_x=\left\lfloor\frac{x}{\delta_{\mathrm L}}\right\rfloor
\quad(x\in[0,1]).
\]
Then \(j_x\delta_{\mathrm L}\le x<(j_x+1)\delta_{\mathrm L}\).
For \(j<j_x\) or \(j>j_x+1\), the distance from \(x\) to
\(j\delta_{\mathrm L}\) is at least \(\delta_{\mathrm L}\), so
\(\phi_j(x)=0\). The two remaining terms satisfy
\[
\begin{aligned}
\phi_{j_x}(x)
 &=\cos\!\left(\frac{\pi(x-j_x\delta_{\mathrm L})}
                         {2\delta_{\mathrm L}}\right),\\
\phi_{j_x+1}(x)
 &=\sin\!\left(\frac{\pi(x-j_x\delta_{\mathrm L})}
                         {2\delta_{\mathrm L}}\right),\\
\phi_{j_x}(x)^2+\phi_{j_x+1}(x)^2&=1.
\end{aligned}
\]
If \(|x-x_0|\ge h_{\mathrm L}\), then \(q(x)=0\).
If \(|x-x_0|<h_{\mathrm L}\) and \(j\notin\mathcal J\), the
open support of \(\phi_j\) fails to intersect the open macro interval;
hence \(|x-j\delta_{\mathrm L}|\ge\delta_{\mathrm L}\) and
\(\phi_j(x)=0\). Consequently, including the macro and lattice
boundaries,
\[
q(x)\sum_{j\in\mathcal J}\phi_j(x)^2=q(x).
\]

For later support counts, define
\[
\mathcal J_x=\{j\in\mathcal J:|x-j\delta_{\mathrm L}|
                                      \le\delta_{\mathrm L}\},
\qquad
\mathcal J_x^\circ=\{j\in\mathcal J:\phi_j(x)\ne0\}.
\]
The closed-support condition places \(j\) between \(j_x-1\) and
\(j_x+1\), while the preceding two-frame argument gives
\[
|\mathcal J_x|\le3,\qquad |\mathcal J_x^\circ|\le2,
\qquad
|F_\lambda(x)|\le\sum_{j\in\mathcal J_x}|\phi_j(x)|\le3,
\qquad |S_\lambda(x)|\le3.
\]

Uniform binary signs are independent, and pairing encodings that differ
in one coordinate gives
\[
E_\lambda\sigma_{\lambda,j}=0,
\qquad
E_\lambda(\sigma_{\lambda,j}\sigma_{\lambda,j'})
=
\begin{cases}
1,&j=j',\\
0,&j\ne j'.
\end{cases}
\]
Expanding the finite frame sum and its square therefore yields
\[
E_\lambda S_\lambda(x)=0,
\qquad
E_\lambda S_\lambda(x)^2
=q(x)\sum_{j\in\mathcal J}\phi_j(x)^2=q(x).
\]
% lean: CausalSmith.Stat.PrivateCateRoughdesign.perturbation_sign_sum
% lean: CausalSmith.Stat.PrivateCateRoughdesign.perturbation_sign_square_sum_exact
% lean: CausalSmith.Stat.PrivateCateRoughdesign.weighted_frame_square_partition

\item \textbf{The centered overlap.}
The prescribed amplitude satisfies
\[
0<t\le\frac1{16384},
\qquad
0\le\sqrt t\le\frac1{128},
\qquad
|\sqrt t\,S_\lambda(x)|\le\frac3{128},
\qquad
0\le tq(x)\le\frac1{16384}.
\]
In particular,
\[
\left|\sqrt t\,S_\lambda(x)\right|<1,
\qquad
\left|-\sqrt t\,S_\lambda(x)\pm tq(x)\right|
\le\frac3{128}+\frac1{16384}<1.
\]
Thus all three Bernoulli parameters in
\cref{obj:def:cosine-family} belong to \([0,1]\).

For each of the four pairs
\((U,V_Y)\in\{-1,1\}^2\), with \(U=2A-1\) and \(V_Y=2Y-1\),
the conditional cell probability divided by its null probability
\(1/4\) is
\[
L_\lambda(x,U,V_Y)
=(1+U\sqrt t\,S_\lambda(x))
 \bigl(1+V_Y(-\sqrt t\,S_\lambda(x)+Utq(x))\bigr).
\]
Expanding in these four cases and using \((\sqrt t)^2=t\) gives
\[
\begin{aligned}
L_\lambda(x,U,V_Y)
={}&1+U\sqrt t\,S_\lambda(x)
 +V_Y\sqrt t\,(tq(x)-1)S_\lambda(x)\\
&+UV_Yt\bigl(q(x)-S_\lambda(x)^2\bigr).
\end{aligned}
\]
Since \(L_\lambda\) is four times a conditional cell probability,
\[
0\le L_\lambda(x,U,V_Y)\le4.
\]
% lean: CausalSmith.Stat.PrivateCateRoughdesign.conditionalLikelihood_fourier

For two sign vectors, define the overlap integrand and its integral by
\[
\begin{aligned}
\mathcal H_{\lambda,\lambda'}(x)
={}&t\bigl(1+(tq(x)-1)^2\bigr)S_\lambda(x)S_{\lambda'}(x)\\
&+t^2\bigl(q(x)-S_\lambda(x)^2\bigr)
       \bigl(q(x)-S_{\lambda'}(x)^2\bigr),
\qquad x\in[0,1],\\
H_{\lambda,\lambda'}&=\int_0^1
                         \mathcal H_{\lambda,\lambda'}(x)\,dx.
\end{aligned}
\]
Expanding the four fair mark cells cancels every product of distinct
members of \(1,U,V_Y,UV_Y\), and gives
\begin{equation}
\label{eq:positive-family-fair-mark-identity}
\frac14\sum_{U\in\{-1,1\}}\sum_{V_Y\in\{-1,1\}}
 L_\lambda(x,U,V_Y)L_{\lambda'}(x,U,V_Y)
=1+\mathcal H_{\lambda,\lambda'}(x).
\end{equation}
Each summand lies in \([0,16]\), so
\[
-1\le\mathcal H_{\lambda,\lambda'}(x)\le15,
\qquad
-1\le H_{\lambda,\lambda'}\le15,
\qquad
|H_{\lambda,\lambda'}|\le15.
\]
All the displayed functions are measurable and bounded.

For fixed \(\lambda'\), the sign moments established above imply
\[
\begin{aligned}
E_\lambda\mathcal H_{\lambda,\lambda'}(x)
={}&t\bigl(1+(tq(x)-1)^2\bigr)
       S_{\lambda'}(x)E_\lambda S_\lambda(x)\\
&+t^2\bigl(q(x)-E_\lambda S_\lambda(x)^2\bigr)
       \bigl(q(x)-S_{\lambda'}(x)^2\bigr)
=0.
\end{aligned}
\]
Commuting the finite sign average with the integral gives
\[
E_\lambda H_{\lambda,\lambda'}=0.
\]
% lean: CausalSmith.Stat.PrivateCateRoughdesign.signOverlap_abs_le_fifteen
% lean: CausalSmith.Stat.PrivateCateRoughdesign.signOverlap_sign_sum

\item \textbf{Coordinate changes and their squared budget.}
Cauchy--Schwarz on the nonzero frames gives
\[
\begin{aligned}
S_\lambda(x)^2
&=q(x)\left(
  \sum_{j\in\mathcal J_x^\circ}\sigma_{\lambda,j}\phi_j(x)
       \right)^2\\
&\le q(x)|\mathcal J_x^\circ|
               \sum_{j\in\mathcal J_x^\circ}\phi_j(x)^2\\
&\le2q(x)\sum_{j\in\mathcal J}\phi_j(x)^2
=2q(x).
\end{aligned}
\]
The same calculation covers an empty nonzero-frame set. Hence
\[
0\le S_\lambda(x)^2\le2q(x),
\qquad
|q(x)-S_\lambda(x)^2|\le q(x),
\qquad
|S_\lambda(x)|\le\frac32.
\]

For \(j\in\mathcal J\) and \(\beta\in\{0,1\}\), define the updated
encoding by
\[
\lambda^{(j,\beta)}_{j'}=
\begin{cases}
\beta,&j'=j,\\
\lambda_{j'},&j'\ne j,
\end{cases}
\qquad j'\in\mathcal J.
\]
Subtracting the two frame sums gives
\[
S_\lambda(x)-S_{\lambda^{(j,\beta)}}(x)
=2g(x)(\lambda_j-\beta)\phi_j(x),
\]
and therefore
\[
|S_\lambda(x)-S_{\lambda^{(j,\beta)}}(x)|
\le2g(x)|\phi_j(x)|\le2.
\]
Both squared perturbations belong to \([0,2q(x)]\), so
\[
\left|S_{\lambda^{(j,\beta)}}(x)^2-S_\lambda(x)^2\right|
\le2q(x)\le2.
\]
Also, \(0\le tq(x)\le1\) implies
\[
0\le1+(tq(x)-1)^2\le2.
\]
The exact difference of the overlap integrands is
\[
\begin{aligned}
\mathcal H_{\lambda,\lambda'}(x)
-\mathcal H_{\lambda^{(j,\beta)},\lambda'}(x)
={}&t\bigl(1+(tq(x)-1)^2\bigr)
   \bigl(S_\lambda(x)-S_{\lambda^{(j,\beta)}}(x)\bigr)
   S_{\lambda'}(x)\\
&+t^2\bigl(S_{\lambda^{(j,\beta)}}(x)^2-S_\lambda(x)^2\bigr)
       \bigl(q(x)-S_{\lambda'}(x)^2\bigr).
\end{aligned}
\]
The absolute values of the two displayed terms are bounded, respectively,
by
\[
t\cdot2\cdot2\cdot\frac32=6t,
\qquad
t^2\cdot2\cdot1=2t^2.
\]
Since \(0\le t\le1\), this proves
\[
\left|\mathcal H_{\lambda,\lambda'}(x)
-\mathcal H_{\lambda^{(j,\beta)},\lambda'}(x)\right|
\le6t+2t^2\le8t.
\]

Define the open frame support within the covariate interval by
\[
\mathcal S_j=\{x\in[0,1]:|x-j\delta_{\mathrm L}|
                                      <\delta_{\mathrm L}\}.
\]
Outside \(\mathcal S_j\), including its endpoints,
\(\phi_j(x)=0\), so the updated and original integrands agree.
Moreover,
\[
\mathcal S_j\subset
(j\delta_{\mathrm L}-\delta_{\mathrm L},
 j\delta_{\mathrm L}+\delta_{\mathrm L}),
\qquad
\operatorname{Leb}(\mathcal S_j)\le2\delta_{\mathrm L}.
\]
It follows that
\[
\begin{aligned}
|H_{\lambda,\lambda'}-H_{\lambda^{(j,\beta)},\lambda'}|
&\le\int_0^1
 \left|\mathcal H_{\lambda,\lambda'}(x)
       -\mathcal H_{\lambda^{(j,\beta)},\lambda'}(x)\right|\,dx\\
&\le8t\,\operatorname{Leb}(\mathcal S_j)
\le16t\delta_{\mathrm L}.
\end{aligned}
\]
% lean: CausalSmith.Stat.PrivateCateRoughdesign.signOverlap_update_abs_le

To count the coordinates, the strict intersection conditions defining
\(\mathcal J\) imply
\[
\frac{x_0-h_{\mathrm L}}{\delta_{\mathrm L}}<j+1,
\qquad
j-1<\frac{x_0+h_{\mathrm L}}{\delta_{\mathrm L}}.
\]
Taking integer bounds yields
\[
\mathcal J\subset
\left[
\left\lfloor\frac{x_0-h_{\mathrm L}}{\delta_{\mathrm L}}\right\rfloor,
\left\lceil\frac{x_0+h_{\mathrm L}}{\delta_{\mathrm L}}\right\rceil
\right]\cap\mathbb Z.
\]
Counting this integer interval and using the floor and ceiling bounds,
\[
\begin{aligned}
|\mathcal J|
&\le
\left\lceil\frac{x_0+h_{\mathrm L}}{\delta_{\mathrm L}}\right\rceil
-\left\lfloor\frac{x_0-h_{\mathrm L}}{\delta_{\mathrm L}}\right\rfloor+1\\
&\le\frac{2h_{\mathrm L}}{\delta_{\mathrm L}}+3
\le\frac{5h_{\mathrm L}}{\delta_{\mathrm L}}.
\end{aligned}
\]
The last inequality uses
\(0<\delta_{\mathrm L}=h_{\mathrm L}^5\le h_{\mathrm L}\).
Thus the squared coordinate-change budget satisfies
\[
\begin{aligned}
\sum_{j\in\mathcal J}(16t\delta_{\mathrm L})^2
&=|\mathcal J|(16t\delta_{\mathrm L})^2\\
&\le\frac{5h_{\mathrm L}}{\delta_{\mathrm L}}
                  (16t\delta_{\mathrm L})^2
=1280t^2h_{\mathrm L}\delta_{\mathrm L}.
\end{aligned}
\]
% lean: CausalSmith.Stat.PrivateCateRoughdesign.sign_coordinate_squared_budget

\item \textbf{The joint exponential average.}
For fixed \(\lambda'\), the function
\(\lambda\mapsto H_{\lambda,\lambda'}\) is a bounded measurable
function of independent fair binary coordinates, has mean zero,
and changes by at most \(16t\delta_{\mathrm L}\) in each coordinate.
Apply the bounded-differences exponential-moment conclusion of
\citep[Theorem~6.2, displayed log-mgf bound in its proof,
p.~167]{BoucheronLugosiMassart2013}, with exponential parameter
\(n\ge0\). It gives
\[
\begin{aligned}
E_\lambda\exp(nH_{\lambda,\lambda'})
&\le
\exp\!\left\{\frac{n^2}{8}
       \sum_{j\in\mathcal J}(16t\delta_{\mathrm L})^2\right\}\\
&\le\exp\!\left\{160n^2t^2h_{\mathrm L}\delta_{\mathrm L}\right\}.
\end{aligned}
\]
Here uniform averaging is the finite product expectation:
\[
E_\lambda\exp(nH_{\lambda,\lambda'})
=\frac1{2^{|\mathcal J|}}
  \sum_{\lambda\in\{0,1\}^{\mathcal J}}\exp(nH_{\lambda,\lambda'}).
\]
Its denominator is positive, including for an empty coordinate set,
when the product law is concentrated on the unique empty encoding.

Let \(\lambda'\) be an independent uniform encoding, and define
\(E_{\lambda,\lambda'}=E_{\lambda'}E_\lambda\).
Averaging the fixed-\(\lambda'\) inequality gives
\[
\begin{aligned}
E_{\lambda,\lambda'}\exp(nH_{\lambda,\lambda'})
&=\frac1{2^{2|\mathcal J|}}
 \sum_{\lambda,\lambda'\in\{0,1\}^{\mathcal J}}
                      \exp(nH_{\lambda,\lambda'})\\
&\le\exp\!\left\{160n^2t^2h_{\mathrm L}\delta_{\mathrm L}\right\}.
\end{aligned}
\]
% lean: CausalSmith.Stat.PrivateCateRoughdesign.signOverlap_exp_average_le
% lean: CausalSmith.Stat.PrivateCateRoughdesign.signOverlap_joint_exp_average_le

\item \textbf{From overlap to chi-square divergence.}
Each conditional mass of \((A,Y(0),Y(1))\) under an alternative is
a product of three Bernoulli masses and is at most \(1\).
The corresponding null mass is \(1/8\).
The covariate law is uniform in both experiments, and the consistency
map is shared. Integrating these conditional inequalities and then
taking the observed marginal gives
\[
P_\lambda\le8P_0,
\qquad
P_\lambda^{\mathrm{obs}}\le8P_0^{\mathrm{obs}},
\qquad
P_\lambda^{\mathrm{obs}}\ll P_0^{\mathrm{obs}}.
\]

Use the observed-record space
\(\mathcal O=[0,1]\times\{0,1\}^2\), and choose the one-record density
version
\[
\mathscr L_\lambda(x,A,Y)
=L_\lambda(x,2A-1,2Y-1)
=\frac{dP_\lambda^{\mathrm{obs}}}{dP_0^{\mathrm{obs}}}(x,A,Y),
\qquad (x,A,Y)\in\mathcal O.
\]
For an ordered dataset
\(D=((X_i,A_i,Y_i))_{i=1}^n\in\mathcal O^n\), define
\[
\mathscr D_n(D)
=E_\lambda\prod_{i=1}^n\mathscr L_\lambda(X_i,A_i,Y_i)
=\frac{dQ_n}{d(P_0^{\mathrm{obs}})^{\otimes n}}(D).
\]
These versions are bounded; in particular,
\(0\le\mathscr D_n\le8^n\).
Finite averaging and product integration therefore apply to all their
first and second moments.

Integrating \cref{eq:positive-family-fair-mark-identity} gives
\[
\int_{\mathcal O}\mathscr L_\lambda\mathscr L_{\lambda'}
                      \,dP_0^{\mathrm{obs}}
=1+H_{\lambda,\lambda'}\ge0.
\]
Since \(\int\mathscr D_n\,d(P_0^{\mathrm{obs}})^{\otimes n}=1\),
expanding the centered square, the finite mixture, and the product
integral gives
\[
\begin{aligned}
1+\chi^2\!\left(Q_n,(P_0^{\mathrm{obs}})^{\otimes n}\right)
&=\int_{\mathcal O^n}\mathscr D_n(D)^2\,
                         d(P_0^{\mathrm{obs}})^{\otimes n}(D)\\
&=E_{\lambda,\lambda'}
 \left(\int_{\mathcal O}
       \mathscr L_\lambda\mathscr L_{\lambda'}\,dP_0^{\mathrm{obs}}
 \right)^n\\
&=E_{\lambda,\lambda'}(1+H_{\lambda,\lambda'})^n.
\end{aligned}
\]
For each pair of encodings, the inequality \(1+H_{\lambda,\lambda'}
\le\exp(H_{\lambda,\lambda'})\), together with
\(1+H_{\lambda,\lambda'}\ge0\), implies
\[
(1+H_{\lambda,\lambda'})^n
\le\exp(nH_{\lambda,\lambda'}).
\]
The empty product and every zeroth power equal \(1\), so these
identities and inequalities include \(n=0\). Combining with the joint
exponential bound and subtracting one proves
\[
\chi^2\!\left(Q_n,(P_0^{\mathrm{obs}})^{\otimes n}\right)
\le
\exp\!\left\{160n^2t^2h_{\mathrm L}\delta_{\mathrm L}\right\}-1.
\]
% lean: CausalSmith.Stat.PrivateCateRoughdesign.signMixture_chiSq_overlap_identity
% lean: CausalSmith.Stat.PrivateCateRoughdesign.signDensityOverlap_eq_one_add_signOverlap
% lean: CausalSmith.Stat.PrivateCateRoughdesign.signMixture_chiSq_le_fourier_exp_average

\item \textbf{Regularity and alternative membership.}
We next verify the smoothness conditions using the two scales of the
construction. For \(\varrho>0\), \(\zeta\in\mathbb R\), and
\(x\in[0,1]\), define the localized cosine
\[
\mathcal C_{\varrho,\zeta}(x)=
\begin{cases}
\cos\!\left(\dfrac{\pi(x-\zeta)}{2\varrho}\right),
   &|x-\zeta|\le\varrho,\\
0,&|x-\zeta|>\varrho.
\end{cases}
\]
It has the clamped representation
\[
\mathcal C_{\varrho,\zeta}(x)
=\cos\!\left(
 \frac{\pi}{2\varrho}
 \max\{-\varrho,\min\{\varrho,x-\zeta\}\}\right).
\]
Indeed, beyond either endpoint the clamped cosine equals
\(\cos(\pm\pi/2)=0\). Clamping is \(1\)-Lipschitz and cosine is
\(1\)-Lipschitz, so
\[
|\mathcal C_{\varrho,\zeta}(x)
 -\mathcal C_{\varrho,\zeta}(x')|
\le\frac{\pi}{2\varrho}|x-x'|.
\]
Also, for two values of absolute value at most \(1\), factoring the
difference of their squares bounds that difference by twice their
distance. Since
\[
\phi_j=\mathcal C_{\delta_{\mathrm L},j\delta_{\mathrm L}},
\qquad
g=\mathcal C_{h_{\mathrm L},x_0}^{\,2},
\qquad q=g^2,
\]
we obtain
\[
\begin{aligned}
|\phi_j(x)-\phi_j(x')|
 &\le\frac{\pi}{2\delta_{\mathrm L}}|x-x'|,\\
|g(x)-g(x')|
 &\le\frac{\pi}{h_{\mathrm L}}|x-x'|,\\
|q(x)-q(x')|
 &\le\frac{2\pi}{h_{\mathrm L}}|x-x'|.
\end{aligned}
\]
% lean: CausalSmith.Stat.PrivateCateRoughdesign.localized_cosine_difference
% lean: CausalSmith.Stat.PrivateCateRoughdesign.bump_difference_bound

Only indices in \(\mathcal J_x\cup\mathcal J_{x'}\) contribute to
\(F_\lambda(x)-F_\lambda(x')\), and that union has at most six
elements. Hence
\[
\begin{aligned}
|F_\lambda(x)-F_\lambda(x')|
&\le\sum_{j\in\mathcal J_x\cup\mathcal J_{x'}}
                       |\phi_j(x)-\phi_j(x')|\\
&\le6\frac{\pi}{2\delta_{\mathrm L}}|x-x'|
=\frac{3\pi}{\delta_{\mathrm L}}|x-x'|.
\end{aligned}
\]
Using \(|F_\lambda(x')|\le3\) and the product identity
\[
S_\lambda(x)-S_\lambda(x')
=g(x)\bigl(F_\lambda(x)-F_\lambda(x')\bigr)
 +\bigl(g(x)-g(x')\bigr)F_\lambda(x'),
\]
we obtain, since \(\delta_{\mathrm L}\le h_{\mathrm L}\) and
\(\pi<4\),
\[
\begin{aligned}
|S_\lambda(x)-S_\lambda(x')|
&\le\left(\frac{3\pi}{\delta_{\mathrm L}}
                  +\frac{3\pi}{h_{\mathrm L}}\right)|x-x'|\\
&\le\frac{6\pi}{\delta_{\mathrm L}}|x-x'|
\le\frac{24}{\delta_{\mathrm L}}|x-x'|.
\end{aligned}
\]
Similarly,
\[
|q(x)-q(x')|
\le\frac{2\pi}{h_{\mathrm L}}|x-x'|
\le\frac{24}{\delta_{\mathrm L}}|x-x'|.
\]
The uniform bounds give
\[
|S_\lambda(x)-S_\lambda(x')|\le6\le24,
\qquad
|q(x)-q(x')|\le1\le24.
\]

If \(x=x'\), both differences vanish. If
\(0<|x-x'|/\delta_{\mathrm L}\le1\), use
\[
\frac{|x-x'|}{\delta_{\mathrm L}}
\le
\left(\frac{|x-x'|}{\delta_{\mathrm L}}\right)^{1/10}
\]
in the Lipschitz bounds. If
\(|x-x'|/\delta_{\mathrm L}>1\), use
\[
1\le
\left(\frac{|x-x'|}{\delta_{\mathrm L}}\right)^{1/10}
\]
in the uniform bounds. Thus, for every \(x,x'\in[0,1]\),
\[
\begin{aligned}
|S_\lambda(x)-S_\lambda(x')|
&\le24\frac{|x-x'|^{1/10}}{\delta_{\mathrm L}^{1/10}},\\
|q(x)-q(x')|
&\le24\frac{|x-x'|^{1/10}}{\delta_{\mathrm L}^{1/10}}.
\end{aligned}
\]
% lean: CausalSmith.Stat.PrivateCateRoughdesign.perturbation_holder_bound
% lean: CausalSmith.Stat.PrivateCateRoughdesign.bump_holder_bound

The scale calibration is
\[
\delta_{\mathrm L}^{1/10}=h_{\mathrm L}^{1/2},
\qquad
\sqrt t=2^{-6}h_{\mathrm L}^{1/2},
\qquad
\frac{\sqrt t}{\delta_{\mathrm L}^{1/10}}=\frac1{64},
\qquad
0\le t\le\sqrt t.
\]
The selected propensity and control mean in
\cref{obj:def:cosine-family} therefore satisfy
\[
\begin{aligned}
|e_{P_\lambda}(x)-e_{P_\lambda}(x')|
&=\frac{\sqrt t}{2}|S_\lambda(x)-S_\lambda(x')|\\
&\le\frac3{16}|x-x'|^{1/10}
\le L|x-x'|^{1/10},
\end{aligned}
\]
and
\[
\begin{aligned}
|\mu_{0,P_\lambda}(x)-\mu_{0,P_\lambda}(x')|
&\le\frac12\left(
 \sqrt t\,|S_\lambda(x)-S_\lambda(x')|
 +t\,|q(x)-q(x')|\right)\\
&\le\frac38|x-x'|^{1/10}
\le L|x-x'|^{1/10},
\end{aligned}
\]
where \cref{obj:synth_4} supplies \(L=3\).
Subtracting the two arm means cancels the shared perturbation:
\[
\tau_{P_\lambda}(x)=tq(x).
\]
Consequently,
\[
\begin{aligned}
|\tau_{P_\lambda}(x)-\tau_{P_\lambda}(x')|
&\le t\frac{2\pi}{h_{\mathrm L}}|x-x'|\\
&=2\pi\,2^{-12}|x-x'|
\le L|x-x'|.
\end{aligned}
\]
These establish
\cref{obj:ass:propensity-holder,obj:ass:control-holder,obj:ass:contrast-lipschitz}.
The amplitude bound also gives
\[
\frac{1-3/128}{2}\le e_{P_\lambda}(x)
\le\frac{1+3/128}{2},
\qquad
\frac14\le e_{P_\lambda}(x)\le\frac34,
\]
as required by \cref{obj:ass:overlap}.
The construction in \cref{obj:def:cosine-family} supplies consistency,
conditional exchangeability, and density \(f_{P_\lambda}=1\), which
satisfy
\cref{obj:ass:consistency,obj:ass:exchangeability,obj:ass:density}.
Every defining condition in \cref{obj:def:complete-model} is therefore
satisfied:
\[
P_\lambda\in\mathcal P\qquad
(\lambda\in\{0,1\}^{\mathcal J}).
\]
% lean: CausalSmith.Stat.PrivateCateRoughdesign.cosineFamily_completeModel

\item \textbf{The null, targets, and one-record cancellation.}
Under \(P_0\), consistency holds by construction, treatment is
conditionally independent of the potential outcomes, and the selected
functions satisfy
\[
f_{P_0}(x)=1,\qquad
e_{P_0}(x)=\mu_{0,P_0}(x)=\frac12,\qquad
\tau_{P_0}(x)=0
\quad(x\in[0,1]).
\]
The density and overlap bounds follow immediately; each smoothness
difference is zero. Thus \cref{obj:def:complete-model} gives
\[
P_0\in\mathcal P,\qquad \theta(P_0)=0.
\]
At the localization center,
\[
g(x_0)=\cos^2(0)=1,\qquad q(x_0)=1,
\qquad
\theta(P_\lambda)=\tau_{P_\lambda}(x_0)=t.
\]

The likelihood expansion established above is the asserted formula.
Its sign average is, pointwise in every covariate and mark pair,
\[
\begin{aligned}
E_\lambda L_\lambda(x,U,V_Y)
={}&1+U\sqrt t\,E_\lambda S_\lambda(x)
 +V_Y\sqrt t\,(tq(x)-1)E_\lambda S_\lambda(x)\\
&+UV_Yt\bigl(q(x)-E_\lambda S_\lambda(x)^2\bigr)
=1.
\end{aligned}
\]
Hence the uniformly averaged conditional probability of each observed
cell is \(1/4\). Integrating against the common uniform covariate law
identifies the measures:
\[
E_\lambda P_\lambda^{\mathrm{obs}}=P_0^{\mathrm{obs}}.
\]
The pointwise frame identities ensure this cancellation also at all
macro and frame boundaries.
% lean: CausalSmith.Stat.PrivateCateRoughdesign.fairNull_completeModel
% lean: CausalSmith.Stat.PrivateCateRoughdesign.fairNull_target
% lean: CausalSmith.Stat.PrivateCateRoughdesign.cosineFamily_target
% lean: CausalSmith.Stat.PrivateCateRoughdesign.oneRecordMixture_eq_fairNull

\item \textbf{Absolute continuity and square integrability.}
Taking finite products of the observed domination inequality and
then averaging gives
\[
(P_\lambda^{\mathrm{obs}})^{\otimes n}
\le8^n(P_0^{\mathrm{obs}})^{\otimes n},
\qquad
Q_n=E_\lambda(P_\lambda^{\mathrm{obs}})^{\otimes n}
\le8^n(P_0^{\mathrm{obs}})^{\otimes n}.
\]
Thus
\[
Q_n\ll(P_0^{\mathrm{obs}})^{\otimes n},
\qquad
0\le\frac{dQ_n}{d(P_0^{\mathrm{obs}})^{\otimes n}}
\le8^n
\quad (P_0^{\mathrm{obs}})^{\otimes n}\text{-almost everywhere}.
\]
The centered squared density is measurable and obeys
\[
0\le
\left(\frac{dQ_n}{d(P_0^{\mathrm{obs}})^{\otimes n}}-1\right)^2
\le(8^n+1)^2.
\]
Since the dominating law is a probability measure,
\[
\int_{\mathcal O^n}
\left(\frac{dQ_n}{d(P_0^{\mathrm{obs}})^{\otimes n}}-1\right)^2
\,d(P_0^{\mathrm{obs}})^{\otimes n}
\le(8^n+1)^2<\infty.
\]
For \(n=0\), both dataset laws are concentrated on the empty dataset
and the density equals \(1\), so the centered integral is zero.
This completes all the asserted certificates.
% lean: CausalSmith.Stat.PrivateCateRoughdesign.signMixture_absolutelyContinuous
% lean: CausalSmith.Stat.PrivateCateRoughdesign.signMixture_square_density_integrable
\end{enumerate}
\end{proof}
